\section{Sviluppo personaggi e obiettivi}

La progettazione dell'applicazione \textbf{Culturia} è stata guidata dalla necessità di soddisfare un'utenza variegata ma accomunata dalla medesima esperienza di fruizione base. Secondo le specifiche di progetto, gli attori principali del sistema non presentano differenze a livello \textit{funzionale}, ma esclusivamente \textit{concettuale}. Essi si interfacceranno con lo stesso applicativo e disporranno dei medesimi strumenti (consultazione della mappa, lettura delle schede informative, navigazione per categorie o ricerca), ma li adopereranno con fini, motivazioni e passioni diametralmente opposte.

Questa scelta architetturale consente di mantenere una piattaforma lineare e inclusiva, delegando alla flessibilità strutturale dei contenuti il compito di assecondare molteplici scenari d'uso. I tre macro-attori cardine del sistema sono: il Turista Esterno, il Turista Interno e l'Amministratore/Impiegato. Nel dettaglio, le analisi hanno imposto un'ulteriore suddivisione del "Turista Interno", generando due flussi distinti.

\subsection{Giustificazione delle scelte e dei Goal}

\textbf{1. Turista Esterno}\\
Questo utente rappresenta il visitatore che non possiede alcuna familiarità, o ne possiede molto scarsa, con il sito culturale e con la città ospitante. Costui ha spesso a disposizione un tempo di visita limitato e necessita di coordinate chiare.
\begin{itemize}
    \item \textbf{Goal primario:} Orientarsi con disinvoltura, identificare e accedere rapidamente alle attrazioni di maggior pregio storico (i classici "must-see"), venendo a contatto con infarinature generali e intuitive.
    \item \textbf{Giustificazione:} Includere il turista esterno è imprescindibile per l'indotto. Considerare questo profilo impone all'app un elevato livello di \textit{learnability}: l'interfaccia non può ammettere frizioni all'apprendimento, consentendo panoramiche fruibili nel limitato tempo concesso per la visita.
\end{itemize}

\textbf{2. Turista Interno (Suddiviso in 2 Personas)}\\
Costituisce il cittadino o abitante della zona. Egli gravita attivamente attorno al sito. L'identificazione del turista interno racchiude due macro-atteggiamenti antitetici:
\begin{itemize}
    \item \textit{Assiduo conoscitore:} Appassionato e avvezzo, padroneggia attivamente le planimetrie locali.
    \begin{itemize}
        \item \textbf{Goal primario:} Esplorare zone di minore popolarità, recuperare curiosità intellettuali o di restauro, scovare dettagli inesplorati sfuggiti nei vecchi sopralluoghi fisici.
    \end{itemize}
    \item \textit{Alla prima visita:} Colui che sa dell'esistenza e del magnetismo del bacino culturale, vi transita attorno o ne legge, ma per apatia o foga cittadina non lo ha mai esplorato.
    \begin{itemize}
        \item \textbf{Goal primario:} Superare il blocco iniziale. Farsi traghettare ad una presa di coscienza sia generale (colmare la base) che disimpegnata con ritmi moderati e svagati.
    \end{itemize}
    \item \textbf{Giustificazione:} Disgiungere l'utenza interna nei due rami è stata una decisione progettuale oculata. Previene, da un lato, l'ostica banalizzazione di \textit{Culturia} (preservando sfide interessanti per i veterani) e al contempo accompagna per mano il locale esordiente non offrendogli interfacce da "snob elitario".
\end{itemize}

\textbf{3. Amministratore / Impiegato}\\
Rappresenta l'operatore preposto al supporto sul posto (steward, custode o centralinista sul campo). Utilizza la \textit{medesima} veste grafica di \textit{Culturia} rivolta alla folla turistica, senza dashboard amministrative.
\begin{itemize}
    \item \textbf{Goal primario:} Interfacciarsi per riscontri documentali logistici e ausiliari. Affiancare e veicolare l'utenza rintracciando dal tablet/smartphone percorsi, uscite d'emergenza o servizi igienici, ed elargire risposte rigorose e fulminee.
    \item \textbf{Giustificazione:} Induce a progettare un'esperienza non rivolta in via esclusiva allo svago. L'impiegato usa l'app quale vero e proprio "attrezzo di servizio": esige correttezza puntuale dei capitoli, nomenclature aderenti allo stato fattuale delle teche, consolidando l'affidabilità del sistema ai massimi gradini.
\end{itemize}
